\documentclass[a4paper,14pt]{extarticle}
\usepackage{cmap}
\usepackage[T2A]{fontenc}
\usepackage[utf8]{inputenc}
\usepackage[russian]{babel}
\usepackage{geometry}
\geometry{left=30mm, right=15mm, top=20mm, bottom=20mm}
\usepackage{indentfirst}
\usepackage{booktabs}
\usepackage{tabularx}
\usepackage{amsmath,amssymb}
\usepackage{xurl}
\usepackage[colorlinks=true,linkcolor=blue,citecolor=blue,urlcolor=blue]{hyperref}

\begin{document}

\begin{center}
    \textbf{\Large Оптимизация поиска базовых чанков в системах дедупликации данных на основе сходства (SBC)}
    
    \vspace{0.5em}
    \textbf{\large Задание 2: Сравнительный анализ подходов, архитектурные выводы и обоснование постановки задачи}
\end{center}

\vspace{1em}

\section*{1. Сравнительный анализ исследованных методов и инструментов}

На основе выполненного обзора литературы проведено качественное и архитектурное сопоставление ключевых подходов к обнаружению сходства блоков и дельта-компрессии. Сравнительная характеристика сведена в табл.~\ref{tab:comparison_analysis}.

\begin{table}[h!]
\small
\centering
\caption{Сравнительный анализ алгоритмов обнаружения сходства и дедупликации}
\label{tab:comparison_analysis}
\begin{tabularx}{\textwidth}{p{3.8cm}ccccX}
\toprule
\textbf{Метод / Алгоритм} & \textbf{Сжатие} & \textbf{CPU} & \textbf{RAM} & \textbf{FP} & \textbf{Ключевое ограничение} \\
\midrule
Классический CDC (LBFS) & Базовое & Низкая & Мин. & Нет & Не находит квази-дубликаты \\
Odess (CDS Sampling) & Высокое & Средняя & Низкая & Средний & Разрастание постингов в индексе \\
Finesse (Feature Locality) & Высокое & Высокая & Средняя & Низкий & Накладные расходы на деление \\
All-Pairs Inverted Index & Среднее & Низкая & Средняя & Низкий & Требует строгого порога \\
\textbf{Marline Palantir} & \textbf{Макс.} & \textbf{Мин.} & \textbf{Низкая} & \textbf{Мин.} & \textbf{Устранено связкой эвристик} \\
\bottomrule
\end{tabularx}
\end{table}

\section*{2. Анализ узких мест существующих подходов}

Анализ теоретических моделей и практических реализаций позволил выявить три фундаментальных архитектурных узких места, ограничивающих производительность систем SBC:

\begin{enumerate}
    \item \textbf{Проблема «тяжёлого хвоста» инвертированного индекса:}
    При индексировании сотен тысяч версий файлов частотные фичи MinHash/Рабина образуют списки постингов длиной в тысячи элементов. Согласно распределению Ципфа, небольшая доля наиболее популярных признаков генерирует более 80\% вычислительной нагрузки при обходе индекса. Время поиска кандидата начинает деградировать квадратично относительно глубины версионной истории.

    \item \textbf{Разрыв между скетч-индексом и энкодером (Size Skew):}
    Существующие алгоритмы скетчирования (Odess, Finesse) работают с эскизами фиксированной длины и полностью абстрагированы от реального физического размера исходных блоков данных. Алгоритм нарезки FastCDC генерирует блоки с дисперсией длины от 100~байт до 64~КБ. Если скетч-индекс находит сходство между чанком длиной 4~КБ и базой длиной 32~КБ, дельта-энкодер GDelta производит полный жадный побайтовый поиск и запускает сжатие Zstd для дельты и оригинала. Однако из-за разницы размеров дельта гарантированно не может дать эффективного коэффициента сжатия и отбрасывается по порогу ложного срабатывания (False Positive, FP). Это приводит к миллионам холостых запусков тяжелого дельта-конвейера.

    \item \textbf{Паразитные дельта-цепочки (Chain Bloat):}
    Классическая схема SBC безусловно привязывает дельтой любой блок, если дельта сжалась хоть немного лучше оригинального блока ($ratio < 0.9$). В результате блоки, у которых дельта составляет 90--95\% от размера базы, сохраняются в виде дельт. Это не даёт ощутимой экономии дискового пространства, но формирует глубокие цепочки восстановления ($C_5 \to C_4 \to C_3 \to C_2 \to C_1 \to \text{Base}$), что катастрофически снижает скорость случайного чтения при декомпрессии.
\end{enumerate}

\section*{3. Обоснование выбора методов, стека технологий и архитектуры}

На основе анализа узких мест сформулированы следующие архитектурные решения:

\begin{itemize}
    \item \textbf{Обоснование языка программирования Rust:}
    Системы вторичного хранения и дедупликации данных функционируют в условиях жестких ограничений по латентности и пропускной способности I/O. Язык Rust обеспечивает:
    \begin{itemize}
        \item Нулевую стоимость абстракций (zero-cost abstractions) и отсутствие сборщика мусора (GC), что исключает паузы при обработке очередей NVMe;
        \item Строгие гарантии безопасности работы с памятью и многопоточностью на этапе компиляции;
        \item Эффективную векторизацию (SIMD) и компиляцию в высокооптимизированный машинный код (-O3, LTO).
    \end{itemize}

    \item \textbf{Обоснование чанкера FastCDC:}
    Использование FastCDC обусловлено его превосходством над классическим Rabin-CDC по скорости сканирования (в 3--10 раз быстрее за счет хеша Gear) и нормализацией длины блока вокруг среднего целевого значения, что снижает дисперсию размеров.

    \item \textbf{Сквозной многоуровневый фильтрующий конвейер:}
    Для преодоления выявленных узких мест спроектирован конвейер взаимодополняющих эвристик:
    \begin{align}
        \text{FastCDC} &\longrightarrow \text{Super-Features} \xrightarrow[\text{rare\_first, max\_df}]{\text{All-Pairs Prefix}} \text{Candidate} \nonumber \\
        &\xrightarrow[\text{Size-Skew}]{\text{Pre-Filter}} \text{GDelta} \xrightarrow[\alpha]{\text{Promote}} \text{Storage}.
    \end{align}
    \begin{enumerate}
        \item \textbf{All-Pairs Prefix Filtering и \texttt{rare\_first}:} обход инвертированного индекса в порядке возрастания документной частоты с досрочным завершением цикла candidate generation на префиксе длины $(m - \min\_overlap + 1)$.
        \item \textbf{Частотная отсечка \texttt{max\_df}:} исключение списков постингов, превышающих пороговое значение длины.
        \item \textbf{Size-Skew Pre-Filter:} мгновенная блокировка вызова энкодера до запуска GDelta, если отношение длин блока и базы превышает $\text{max\_size\_skew}$ (например $> 2.0..2.5$).
        \item \textbf{Min-Chunk Pre-Filter:} исключение LSH-поиска для микро-чанков ($< 1$~КБ).
        \item \textbf{Порог $\alpha$ (Promote to Base):} принудительное сохранение блока как новой самостоятельной базы, если размер дельты превышает долю $\alpha$ от размера базы ($|Delta| \ge \alpha \cdot |Base|$), что гарантированно обрывает паразитные цепочки дельт.
    \end{enumerate}
\end{itemize}

\section*{4. Скорректированная постановка цели и декомпозиция задач}

Проведённый сравнительный анализ обосновывает корректировку формальной цели и задач исследования:

\textbf{Скорректированная цель:} Разработка, математическое обоснование и экспериментальная верификация комплекса алгоритмических эвристик для скетч-индекса и дельта-конвейера платформы \texttt{marline}, обеспечивающего рост пропускной способности дедупликации более чем на 40--50\% при сохранении степени сжатия в пределах 0.5--1.5\% от абсолютного теоретического предела полного перебора.

\textbf{Скорректированная декомпозиция задач:}
\begin{enumerate}
    \item \textbf{Задача 1. Эвристики фильтрации инвертированного скетч-индекса:}
    Аналитический вывод порогов Жаккара $\min\_overlap = \lceil \frac{2nt}{1+t} \rceil$ и точек бифуркации; реализация обхода \texttt{rare\_first}, префиксной фильтрации All-Pairs, частотной отсечки \texttt{max\_df} и динамической метрики неудачных верификаций \texttt{fp\_metric\_threshold}.

    \item \textbf{Задача 2. Пре-фильтрация кандидатов дельта-конвейера (Size-Skew и Min-Chunk):}
    Разработка механизмов блокировки вызова дельта-энкодера при размерном перекосе блоков ($\text{skew} > \text{max\_size\_skew}$) и обхода LSH-поиска для микро-чанков ($|C| < \text{min\_chunk\_size}$).

    \item \textbf{Задача 3. Динамическое управление жизненным циклом баз (порог $\alpha$):}
    Реализация алгоритма разрыва длинных дельта-цепочек через объявление блока новой базой при неэффективной дельте ($|Delta| \ge \alpha \cdot |Base|$).

    \item \textbf{Задача 4. Экспериментальное исследование и кросс-датасетная валидация:}
    Проведение сравнительного бенчмаркинга на каноническом датасете Linux Kernel 3.5.6--3.7.6 (1.31~ГБ) и независимом валидационном датасете Linux Kernel 3.4.5--3.4.7 (1.26~ГБ), построение карт Парето и формирование эксплуатационных профилей.
\end{enumerate}

\end{document}
